\documentclass[11pt,a4paper]{article}
\usepackage[margin=1in]{geometry}
\usepackage{amsmath, mathtools, amssymb, amsthm}
\usepackage{graphicx}
\usepackage{booktabs}
\usepackage{xcolor}
\usepackage[colorlinks=true, linkcolor=blue, citecolor=blue, urlcolor=blue]{hyperref}

\theoremstyle{plain}
\newtheorem{theorem}{Theorem}
\newtheorem{lemma}[theorem]{Lemma}
\newtheorem{proposition}[theorem]{Proposition}
\newtheorem{corollary}[theorem]{Corollary}

\theoremstyle{definition}
\newtheorem{definition}[theorem]{Definition}
\newtheorem{remark}[theorem]{Remark}

% Custom todo box for in-progress items
\newcommand{\todo}[1]{\vspace{0.3em}\noindent\fbox{\parbox{0.95\textwidth}{\textit{\textbf{Draft note:} #1}}}\vspace{0.3em}}

% ============================================================
% TITLE
% ============================================================

\title{\textbf{Paper Title:\\
  Subtitle If Present}}
\author{Stefanos Drakos\\
\small{AGEL AI I.K.E., Rhodes, Greece}\\
\small{ORCID: 0000-0001-7417-2444}}
\date{\today}

\begin{document}
\maketitle

% ============================================================
% ABSTRACT
% ============================================================

\begin{abstract}
One-paragraph summary, roughly 200--300 words. State the problem, the
approach, the key contribution, and one headline result. End with a
sentence locating the paper relative to prior work (e.g., ``These
results extend the framework of [X] by introducing [Y]'').
\end{abstract}

% ============================================================
% 1. INTRODUCTION
% ============================================================

\section{Introduction}
\label{sec:intro}

Text of introduction, typically 3--5 paragraphs:

\paragraph{Paragraph 1: Problem.} State the applied problem and
why it matters. Keep it concrete (``personalized therapy scheduling
for tumor growth'', not ``decision-making under uncertainty'').

\paragraph{Paragraph 2: Prior work.} Locate the paper relative to
the two or three most relevant lines of research. Cite narrowly,
not exhaustively --- this is not a review paper.

\paragraph{Paragraph 3: Contribution.} State the specific
contribution in bullet form if helpful. Be concrete about what's
new versus what extends existing work.

\paragraph{Paragraph 4: Outline.} Brief section-by-section preview.

% ============================================================
% 2. MODEL AND NOTATION
% ============================================================

\section{Model and Notation}
\label{sec:model}

Introduce the SDE/FPE/model. Define every symbol that appears later.

\begin{equation}
\frac{\partial P}{\partial t} = -\frac{\partial}{\partial y}(\mu(y,t) P) + \frac{1}{2}\frac{\partial^2}{\partial y^2}(D(t) P)
\label{eq:fpe}
\end{equation}

% ============================================================
% 3. MAIN RESULT OR METHODS
% ============================================================

\section{Main Result}
\label{sec:main}

\begin{theorem}[Descriptive name]
\label{thm:main}
Let [hypothesis]. Then [conclusion].
\end{theorem}

\begin{proof}
See Appendix~\ref{app:proof}.
\end{proof}

% ============================================================
% 4. NUMERICAL EXPERIMENTS
% ============================================================

\section{Numerical Experiments}
\label{sec:experiments}

Report 3--5 experiments testing the main claim. Each experiment:
\begin{itemize}
\item Setup: parameters, patient population, discretization
\item Metrics: what's measured, expected values under the theorem
\item Results: table or figure with computed values vs predicted
\item Interpretation: one sentence summary per experiment
\end{itemize}

% ============================================================
% 5. DISCUSSION
% ============================================================

\section{Discussion}
\label{sec:discussion}

\subsection{Scope and limitations}

Name explicitly what this paper does not cover.

\subsection{Connections to related work}

Compare to 2--3 closest prior works. Be specific about where
approaches agree and where they differ.

\subsection{Future directions}

Concrete next steps, ideally 2--4 enumerated items.

% ============================================================
% 6. SOFTWARE
% ============================================================

\section{Software}
\label{sec:software}

Brief description of companion package:
\begin{itemize}
\item Language and dependencies
\item Public URL (GitHub, Zenodo DOI)
\item Reproducibility statement (all figures in this paper reproduced
  by running \texttt{scripts/reproduce\_all.py})
\end{itemize}

% ============================================================
% REFERENCES
% ============================================================

\section*{References}

\begin{list}{}{\leftmargin 2em \itemindent -2em \itemsep 4pt}

\item Author, A.\ B., \& Author, C.\ D.\ (YYYY). Title of the paper
in plain text. \emph{Journal Name} volume(issue):pages.
\href{https://doi.org/DOI}{doi:DOI}

\end{list}

% ============================================================
% APPENDIX
% ============================================================

\appendix

\section{Proof of Theorem~\ref{thm:main}}
\label{app:proof}

\textit{\textbf{Status of this appendix:} brief summary of what the
proof covers and what it defers.}

\subsection{Preliminary definitions}

\subsection{Core lemmas}

\begin{lemma}[Name]
\label{lem:support}
Statement.
\end{lemma}

\begin{proof}
Proof body.
\end{proof}

\subsection{Proof of main theorem}

\begin{proof}[Proof of Theorem~\ref{thm:main}]
Proof body using the lemmas.
\end{proof}

\subsection{Scope remark}

\begin{remark}[Scope of the theorem]
Theorem~\ref{thm:main} does not claim:
\begin{enumerate}
\item [specific limitation 1];
\item [specific limitation 2].
\end{enumerate}
\end{remark}

\end{document}
